	\documentclass[14pt]{extreport}
\usepackage{extsizes}
	\usepackage[frenchb]{babel}
	\usepackage[utf8]{inputenc}  
	\usepackage[T1]{fontenc}
	\usepackage{amssymb}
	\usepackage[mathscr]{euscript}
	\usepackage{stmaryrd}
	\usepackage{amsmath}
	\usepackage{tikz}
	\usepackage[all,cmtip]{xy}
	\usepackage{amsthm}
	\usepackage{varioref}
	\usepackage[ margin=1in]{geometry}
	\geometry{a4paper}
	\usepackage{lmodern}
	\usepackage{hyperref}
	\usepackage{array}
	\usepackage{float}
	\usepackage{easytable}
	 \usepackage{fancyhdr}\usepackage{longtable}
	 \usetikzlibrary{shapes.misc}
	 \newcommand\ang[1]{$#1{}^o$}
\newlength{\taillecellule}
\setlength{\taillecellule}{2cm}
\newcolumntype{C}{@{}>{\centering\arraybackslash}p{\taillecellule}@{}}

\usepackage{pstricks,multido}
\usepackage{arrayjob}
\usepackage{calc,xlop}
\tikzset{cross/.style={cross out, draw=black, minimum size=2*(#1-\pgflinewidth), inner sep=0pt, outer sep=0pt},
%default radius will be 1pt. 
cross/.default={1pt}}

	\pagestyle{fancy}
	\theoremstyle{plain}
	\fancyfoot[C]{\empty} 
	\fancyhead[L]{Interrogation chapitre 6}
	\fancyhead[R]{9 avril 2025}
	
	
	\title{Interrogation chapitre 6}
	\date{}
	\begin{document}



\subsection*{Exercice 1 : Calculs (8 pts)}  

Effectuer sur votre copie, en détaillant les étapes les calculs suivants : 
\begin{enumerate}
\item $ 10 - 15 = $ 
\item $ -3 - (-8)= $ 
\item $ -6 - 1 = $

\item $9- 5 \times 3 + 1 =$
\item $ 14 - 2 \times 5 = $
\item $ 12 - (8 - 4) + ( 5 - 1) = $ 

\item $-1 - (-4) + 2 - 5 + (-1) =$
\item $ 11 - [ 12 - (4 - (-1)) + 2]=$
\end{enumerate}

\subsection*{Exercice 2 : Angles alternes-internes(6 pts)}
On fixe les notations de la figure ci-dessus :\\ 
\includegraphics[scale=.5]{Fig1}\\
\begin{enumerate}
\item On donne $\widehat{ABC} = 70^o$, et $(AB)//(CD)$. 
Que vaut l'angle $\widehat{BCD}$. Rédiger la réponse. 
\item On donne $\widehat{ABC} = 60^o$ et $\widehat{BCE}= 120^o$. Les droites  $(AB)$ et $(CD)$ sont-elles parallèles ? Rédiger la réponse.
\end{enumerate}
	
	
\subsection*{Exercice 3 - Problème (6 pts)}
On considère la figure suivante : \\
        \includegraphics[scale=1]{Fig2}
    \begin{enumerate}
    \item Montrer que les droites $(BG)$ et $(CF)$ sont parallèles. Justifier votre réponse. (Indication : utiliser les angles formés avec $(AD)$.)
    \item Calculer l'angle $\widehat{JKL}$. Justifier votre réponse. 
    \item Déterminer les quatre angles du quadrilatère $IJKL$. Que vaut leur somme ? 
    \end{enumerate}
 
    
    
    
    
    \newpage


\subsection*{Exercice 1 : Calculs (8 pts)}  

Effectuer sur votre copie, en détaillant les étapes les calculs suivants : 
\begin{enumerate}
\item $ 10 - 13 = $ 
\item $ -7 - (-8)= $ 
\item $ -3 - 2 = $

\item $11 - 5 \times 3 + 1 =$
\item $ 18 - 2 \times 5 = $
\item $ 12 - (8 - 3) + ( 5 - 1) = $ 

\item $-1 - (-4) + 3 - 5 + (-1) =$
\item $ 11 - [ 12 - (4 - (-1)) + 5]=$
\end{enumerate}

\subsection*{Exercice 2 : Angles alternes-internes(6 pts)}
On fixe les notations de la figure ci-dessus :\\ 
\includegraphics[scale=.5]{Fig1}\\
\begin{enumerate}
\item On donne $\widehat{ABC} = 50^o$, et $(AB)//(CD)$. 
Que vaut l'angle $\widehat{BCD}$. Rédiger la réponse. 
\item On donne $\widehat{ABC} = 60^o$ et $\widehat{BCE}= 120^o$. Les droites  $(AB)$ et $(CD)$ sont-elles parallèles ? Rédiger la réponse.
\end{enumerate}
	
	
\subsection*{Exercice 3 - Problème (6 pts)}
On considère la figure suivante : \\
        \includegraphics[scale=1]{Fig2}
    \begin{enumerate}
    \item Montrer que les droites $(BG)$ et $(CF)$ sont parallèles. Justifier votre réponse. (Indication : utiliser les angles formés avec $(AD)$.)
    \item Calculer l'angle $\widehat{JKL}$. Justifier votre réponse. 
    \item Déterminer les quatre angles du quadrilatère $IJKL$. Que vaut leur somme ? 
    \end{enumerate}
 
\end{document}